% =============================================================
% 07_security.tex — セキュリティとAI
% =============================================================
\chapter{セキュリティとAI}

\chaptermeta{%
  \item CIA三要素を具体例と共に説明できる
  \item 認証の3要素を区別できる
  \item 公開鍵暗号とデジタル署名の仕組みを理解する
  \item 主要な脅威と対策を整理できる
  \item AI/ML/DLの関係と基本用語を説明できる
}{90分}{3}{第5--6章}

テクノロジ系で最も出題数が多いのがセキュリティ分野です。
近年はAI関連の出題も増加傾向にあります。
\textbf{暗号方式の仕組み}と\textbf{攻撃手法の名前}を確実に押さえましょう。

% =============================================================
\section{CIA三要素}
% =============================================================

\begin{teigi}[情報セキュリティのCIA]
\begin{itemize}
  \item \termtag{機密性}（Confidentiality）：許可された人だけがアクセスできること
  \item \termtag{完全性}（Integrity）：情報が改ざん・破壊されていないこと
  \item \termtag{可用性}（Availability）：必要なときに利用できること
\end{itemize}
この3つのバランスを保つことが情報セキュリティの目的。
\end{teigi}

\begin{hikaku}[CIA三要素の具体例]
\comptable{要素 & 脅威の例 & 対策の例}{
機密性 & 不正アクセス・盗聴 & 暗号化・アクセス制御 \\
完全性 & データ改ざん・ウイルス & ハッシュ値検証・デジタル署名 \\
可用性 & DoS攻撃・機器故障 & 冗長化・バックアップ
}
\end{hikaku}

% =============================================================
\section{認証の3要素}
% =============================================================

\begin{teigi}[認証の3要素]
\begin{itemize}
  \item \termtag{知識情報}：本人だけが知っていること（パスワード、PINコード、秘密の質問）
  \item \termtag{所持情報}：本人だけが持っているもの（ICカード、スマートフォン、ワンタイムパスワード生成器）
  \item \termtag{生体情報}：本人の身体的特徴（指紋、顔、虹彩、静脈パターン）
\end{itemize}
\termtag{二要素認証}は上記3種から\textbf{異なる2種}を組み合わせること。
パスワード＋秘密の質問は両方「知識」なので二要素認証にはならない。
\end{teigi}

\begin{pitfall}[「二段階認証」＝「二要素認証」ではない]
\begin{itemize}
  \item 二段階認証：認証を2回行うこと（同じ要素でもよい）
  \item 二要素認証：異なる種類の要素を2つ組み合わせること
  \item パスワード＋SMS認証 → 知識＋所持 → 二要素認証
  \item パスワード＋秘密の質問 → 知識＋知識 → 二段階だが一要素
\end{itemize}
\end{pitfall}

% =============================================================
\section{公開鍵暗号とデジタル署名}
% =============================================================

\begin{hikaku}[共通鍵暗号と公開鍵暗号]
\comptable{項目 & 共通鍵暗号 & 公開鍵暗号}{
鍵の数 & 送受信者で同じ1つの鍵 & 公開鍵と秘密鍵のペア \\
速度 & 高速 & 低速 \\
鍵配送問題 & あり（安全に鍵を渡す必要） & なし（公開鍵は公開してOK） \\
代表例 & AES & RSA
}
\end{hikaku}

\begin{teigi}[デジタル署名の仕組み]
\begin{enumerate}
  \item 送信者がメッセージの\textbf{ハッシュ値}を計算
  \item 送信者の\textbf{秘密鍵}でハッシュ値を暗号化 → 署名
  \item 受信者は送信者の\textbf{公開鍵}で署名を復号
  \item 復号したハッシュ値と、自分で計算したハッシュ値を比較
  \item 一致すれば\textbf{改ざんなし}かつ\textbf{本人が送信}と確認
\end{enumerate}
公開鍵暗号の「逆使い」——暗号化は受信者の公開鍵、署名は送信者の秘密鍵。
\end{teigi}

% --- 例題 ---
\begin{reidai}{公開鍵暗号 vs デジタル署名}
公開鍵暗号方式による暗号化通信とデジタル署名について、正しい組合せはどれか。

\begin{center}
\small
\begin{tabular}{lll}
\toprule
 & 暗号化通信 & デジタル署名 \\
\midrule
ア & 受信者の公開鍵で暗号化 & 送信者の公開鍵で署名 \\
イ & 受信者の公開鍵で暗号化 & 送信者の秘密鍵で署名 \\
ウ & 送信者の秘密鍵で暗号化 & 受信者の公開鍵で署名 \\
エ & 送信者の公開鍵で暗号化 & 受信者の秘密鍵で署名 \\
\bottomrule
\end{tabular}
\end{center}

\sentakushi{ア}{イ}{ウ}{エ}
\end{reidai}

\begin{shishin}
暗号化通信：「相手しか復号できない」→ 受信者の公開鍵で暗号化、受信者の秘密鍵で復号。
デジタル署名：「本人が送った証拠」→ 送信者の秘密鍵で署名、送信者の公開鍵で検証。
\end{shishin}

\begin{kaitou}
\textbf{正解：イ}

暗号化通信 = 受信者の公開鍵で暗号化。デジタル署名 = 送信者の秘密鍵で署名。
「暗号は相手の公開鍵、署名は自分の秘密鍵」と覚える。
\end{kaitou}

\begin{minuki}[見抜き方]
\begin{itemize}
  \item 暗号化：「受信者の公開鍵」→ 受信者だけが復号できる
  \item デジタル署名：「送信者の秘密鍵」→ 送信者本人の証明
  \item 「暗号は受信者、署名は送信者」——主役が入れ替わる
\end{itemize}
\end{minuki}

% =============================================================
\section{脅威と対策}
% =============================================================

\begin{hikaku}[主要なマルウェアと攻撃手法]
\comptable{種類 & 特徴 & 対策}{
ランサムウェア & データを暗号化して身代金を要求 & バックアップ・OS更新 \\
フィッシング & 偽サイトで情報を騙し取る & URL確認・二要素認証 \\
標的型攻撃 & 特定組織を狙った巧妙なメール & 訓練・メールフィルタ \\
DoS/DDoS & 大量通信でサービスを停止させる & CDN・IP制限 \\
SQLインジェクション & 入力フォームからDB操作 & 入力値の無害化 \\
クロスサイトスクリプティング & Webページに悪意のスクリプト注入 & エスケープ処理 \\
ソーシャルエンジニアリング & 人間の心理を利用して情報を入手 & セキュリティ教育
}
\end{hikaku}

\begin{pitfall}[ランサムウェアの対処で「身代金を払う」を選びがち]
\begin{itemize}
  \item 身代金を払ってもデータが復旧する保証はない
  \item 正しい対策は\textbf{日頃からバックアップ}を取ること
  \item 感染時は即座にネットワークから切り離す
\end{itemize}
\end{pitfall}

% =============================================================
\section{信頼性設計}
% =============================================================

\begin{hikaku}[信頼性設計の比較]
\comptable{設計思想 & 方針 & 具体例}{
フォールトトレランス & 故障しても動き続ける & 冗長化・RAID \\
フェールセーフ & 故障時に安全側に倒す & 踏切の遮断機（故障→閉じる） \\
フェールソフト & 故障時に機能を縮退して継続 & 空調故障→換気だけ継続 \\
フールプルーフ & 人間のミスを防ぐ設計 & コンセントの形状・確認ダイアログ
}
\end{hikaku}

% =============================================================
\section{AI概論}
% =============================================================

\begin{teigi}[AI / ML / DLの関係]
\begin{itemize}
  \item \termtag{AI}（人工知能）：人間の知能を模倣する技術の総称
  \item \termtag{ML}（機械学習）：AIの一分野。データからパターンを自動学習
  \item \termtag{DL}（深層学習/ディープラーニング）：MLの一手法。多層ニューラルネットワーク
  \item 関係：AI $\supset$ ML $\supset$ DL（入れ子構造）
\end{itemize}
\end{teigi}

\begin{hikaku}[機械学習の3分類]
\comptable{学習方式 & 特徴 & 例}{
教師あり学習 & 正解ラベル付きデータで学習 & 画像分類・スパム判定 \\
教師なし学習 & 正解なしでパターンを発見 & クラスタリング・次元削減 \\
強化学習 & 報酬を最大化する行動を学習 & ゲームAI・ロボット制御
}
\end{hikaku}

\begin{teigi}[AI関連の重要用語]
\begin{itemize}
  \item \termtag{AIバイアス}：学習データの偏りにより、AIが不公平な判断をすること
  \item \termtag{説明可能AI（XAI）}：AIの判断根拠を人間が理解できるようにする技術
  \item \termtag{生成AI}：テキスト・画像などのコンテンツを生成するAI（例：ChatGPT）
  \item \termtag{過学習}：訓練データに適合しすぎて、未知データへの予測精度が下がる現象
\end{itemize}
\end{teigi}

% =============================================================
\section{過去問ドリル}
% =============================================================

\shoumatsuhead{過去問ドリル：セキュリティとAI}

\begin{itpmon}[\sourcebadge{R7}\enspace\diffbadge{標}]{%
JIS Q 27017に基づく、クラウドサービス利用におけるセキュリティ管理策の説明として、最も適切なものはどれか。}
\sentakushi
  {クラウドサービス利用者が独自にセキュリティ基準を策定する指針}
  {クラウドサービスに関する情報セキュリティ管理策の実践規範}
  {オンプレミス環境のみを対象とした情報セキュリティ規格}
  {個人情報保護に特化したクラウドセキュリティのガイドライン}
\end{itpmon}
\begin{kaisetsu}{イ}
\textbf{ISO/IEC 27017}（JIS Q 27017）はクラウドサービスに関する情報セキュリティ管理策の実践規範。ISO 27001/27002を補完するクラウド固有のガイドライン。
\end{kaisetsu}

\begin{itpmon}[\sourcebadge{R7}\enspace\diffbadge{標}]{%
AIが学習データの偏りによって、特定の属性を持つ人々に対して不公平な判断を行ってしまう問題を何というか。}
\sentakushi
  {過学習}
  {AIバイアス}
  {シンギュラリティ}
  {チューリングテスト}
\end{itpmon}
\begin{kaisetsu}{イ}
\textbf{AIバイアス}は学習データの偏り（性別・人種等）によりAIが不公平な判断をすること。アは訓練データへの過剰適合、ウはAIが人間の知能を超える仮説上の時点、エはAIの知能を判定するテスト。
\end{kaisetsu}

\begin{itpmon}[\sourcebadge{original}\enspace\diffbadge{標}]{%
ランサムウェアに関する説明として、最も適切なものはどれか。}
\sentakushi
  {Webサイトを閲覧しただけで感染し、広告を大量に表示するソフトウェア}
  {コンピュータ内のファイルを暗号化し、復号の見返りに金銭を要求するマルウェア}
  {キーボードの入力内容を記録して外部に送信するソフトウェア}
  {正常なソフトウェアを装ってインストールさせ、裏で不正な動作を行うマルウェア}
\end{itpmon}
\begin{kaisetsu}{イ}
\textbf{ランサムウェア}（ransom=身代金）はファイルを暗号化し、復号と引き換えに身代金を要求する。アはアドウェア、ウはキーロガー、エはトロイの木馬。
\end{kaisetsu}

% =============================================================
\section{タイムアタック}
% =============================================================

\begin{timeattack}{セキュリティ・AIクイックチェック}{6分}{5問中4問}

\begin{itpmon}[\diffbadge{基}]{情報セキュリティのCIA三要素に含まれないものはどれか。}
\sentakushi{機密性}{完全性}{可用性}{効率性}
\end{itpmon}
\begin{kaisetsu}{エ}
CIAは機密性（C）・完全性（I）・可用性（A）。\textbf{効率性}はCIAに含まれない。
\end{kaisetsu}

\begin{itpmon}[\diffbadge{標}]{デジタル署名で、送信者が署名に使うのはどの鍵か。}
\sentakushi{受信者の公開鍵}{受信者の秘密鍵}{送信者の公開鍵}{送信者の秘密鍵}
\end{itpmon}
\begin{kaisetsu}{エ}
デジタル署名は\textbf{送信者の秘密鍵}で署名し、送信者の公開鍵で検証する。
\end{kaisetsu}

\begin{itpmon}[\diffbadge{基}]{フェールセーフの例として最も適切なものはどれか。}
\sentakushi
  {システムの一部が故障しても全体は動き続ける}
  {故障時に安全な状態で停止する}
  {人間がミスをしても危険な動作にならない設計}
  {機能を縮退して運転を継続する}
\end{itpmon}
\begin{kaisetsu}{イ}
\textbf{フェールセーフ}は故障時に安全側に倒す設計。アはフォールトトレランス、ウはフールプルーフ、エはフェールソフト。
\end{kaisetsu}

\begin{itpmon}[\diffbadge{標}]{機械学習で、正解ラベル付きデータを使って学習する方式はどれか。}
\sentakushi{教師あり学習}{教師なし学習}{強化学習}{転移学習}
\end{itpmon}
\begin{kaisetsu}{ア}
\textbf{教師あり学習}は正解ラベル付きデータで学習。イは正解なし、ウは報酬最大化、エは別タスクで学習した知識を転用する手法。
\end{kaisetsu}

\begin{itpmon}[\diffbadge{標}]{SQLインジェクション攻撃の対策として最も適切なものはどれか。}
\sentakushi
  {強力なパスワードを設定する}
  {入力値の無害化（サニタイジング）を行う}
  {ファイアウォールを設置する}
  {通信を暗号化する}
\end{itpmon}
\begin{kaisetsu}{イ}
SQLインジェクションはWebフォームの入力値にSQL文を混入させる攻撃。対策は\textbf{入力値の無害化}（プレースホルダやエスケープ処理）。
\end{kaisetsu}

\end{timeattack}

% =============================================================
\section{よくあるミスと到達確認}
% =============================================================

\begin{yokumiss}[この章でよくあるミス]
\begin{enumerate}[leftmargin=1.6em, itemsep=5pt]
  \item \textbf{公開鍵暗号で「送信者の公開鍵で暗号化」と答える}——暗号化は受信者の公開鍵。
  \item \textbf{デジタル署名で「公開鍵で署名」と答える}——署名は送信者の秘密鍵。
  \item \textbf{フェールセーフとフールプルーフを混同する}——セーフは「故障→安全停止」、プルーフは「人為ミス防止」。
  \item \textbf{二段階認証を二要素認証と同一視する}——異なる種類の要素を使うかどうかが違い。
  \item \textbf{AIバイアスと過学習を混同する}——バイアスはデータの偏り、過学習は訓練データへの過剰適合。
\end{enumerate}
\end{yokumiss}

\begin{tatsaku}
  \item CIA三要素を具体例付きで説明できる
  \item 認証の3要素をそれぞれ2つ以上例示できる
  \item 公開鍵暗号とデジタル署名の鍵の使い分けを即答できる
  \item ランサムウェア・フィッシング・標的型攻撃を区別できる
  \item フォールトトレランス/フェールセーフ/フェールソフト/フールプルーフを区別できる
  \item AI $\supset$ ML $\supset$ DLの関係を説明できる
  \item 教師あり/教師なし/強化学習を区別できる
  \item AIバイアスと説明可能AIの意味を説明できる
\end{tatsaku}

\nextchapter{%
分野横断の整理と試験戦略に進みます。全章の知識をつなぎ、
ひっかけの見抜き方、本試験の時間配分を攻略します。
}
